\documentclass[11pt,a4paper]{article}
\usepackage[T1]{fontenc}
\usepackage{amsmath,amssymb,amsthm,mathtools}
\usepackage{geometry}
\usepackage[hidelinks]{hyperref}
\geometry{margin=1in}

\newtheorem{theorem}{Theorem}[section]
\newtheorem{lemma}[theorem]{Lemma}
\newtheorem{proposition}[theorem]{Proposition}
\newtheorem{corollary}[theorem]{Corollary}
\theoremstyle{definition}
\newtheorem{definition}[theorem]{Definition}
\newtheorem{example}[theorem]{Example}

\author{Part II}
\date{Recorded from the lecturer in 2026-2027.}

\title{Principles of Statistics}
\begin{document}
\maketitle
\thispagestyle{empty}
\section*{Principles of Statistics}

\subsection*{0. INTRODUCTION}

Let $(\Omega,\mathcal{F},\mathbb{P})$ be a probability space, and
$X:\Omega\to\mathbb{R}$ a random variable with distribution function
\[
\begin{aligned}
F(t)
&=
\mathbb{P}\bigl(\{\omega:X(\omega)\leq t\}\bigr) \\
&=
\begin{cases}
\displaystyle \int_{-\infty}^{t} f(x)\,dx,
& \text{if $X$ is continuous with pdf $f$}, \\[1ex]
\displaystyle \sum_{x\leq t} f(x),
& \text{if $X$ is discrete with pmf $f$}.
\end{cases}
\end{aligned}
\]

We denote $P=\mathbb{P}\circ X^{-1}$, the law or ``distribution'' of $X$.

We observe $X_1,\ldots,X_n$, which are iid copies of $X$, where $n$ is
the \underline{sample size}.

\paragraph{Def 1.}
A \underline{statistical model} is a collection
$\{f_\theta:\theta\in\Theta\}$ of pdfs or pmfs, or
$\{P_\theta:\theta\in\Theta\}$ of laws, indexed by a
\underline{parameter space} $\Theta$.

\paragraph{Ex.}
\begin{itemize}
    \item $N(\theta,1)$, \qquad $\Theta=\mathbb{R}$.
    \item $N(\mu,\sigma^2)$, \qquad
    $\theta=(\mu,\sigma^2)$, \quad
    $\Theta=\mathbb{R}\times(0,\infty)$.
    \item $\operatorname{Poi}(\theta)$, \qquad
    $\Theta=[0,\infty)$.
\end{itemize}

\paragraph{Def 2.}
We say the model is \underline{correctly specified} when
$\exists\,\theta_0\in\Theta$ such that $P=P_{\theta_0}$.

We will refer to the ``true parameter'' as $\theta_0$.

\subsection*{Goals of Statistics}

\begin{enumerate}
    \item \underline{Estimation}:
    We define an \underline{estimator}
    \[
    \widehat{\theta}_n
    =
    \widehat{\theta}_n(X_1,\ldots,X_n),
    \]
    which we hope is close to $\theta_0$ when
    $X_i\overset{\mathrm{iid}}{\sim}P_{\theta_0}$,
    for all $\theta_0\in\Theta$, when $n$ is large.

    \item \underline{Testing}:
    We have a null hypothesis
    \[
    H_0:\theta\in\Theta_0\subseteq\Theta,
    \]
    and an alternative hypothesis
    \[
    H_1:\theta\in\Theta_1\subseteq\Theta.
    \]
    We define a \underline{test},
    \[
    \psi_n:=\psi_n(X_1,\ldots,X_n),
    \]
    such that $\psi_n=0$ if $H_0$ is true, and
    $\psi_n=1$ if $H_1$ is true, \underline{w.h.p.}

    \item \underline{Uncertainty Quantification}
    \begin{itemize}
        \item Define a \underline{confidence set}
        $C_n(X_1,\ldots,X_n;\alpha)$, such that
        \[
        P_\theta(\theta\in C_n)
        \overset{(\geq)}{=}
        1-\alpha
        \qquad \forall\,\theta\in\Theta,
        \]
        where $\alpha$ is the \underline{significance level}.

        \item Define a \underline{Bayesian credible set},
        a region of high posterior probability.
    \end{itemize}
\end{enumerate}

\section*{I. LIKELIHOOD PRINCIPLE}

\paragraph{Def 3.}
The \underline{likelihood function}
\[
L_n(\theta)=\prod_{i=1}^{n}f_\theta(X_i).
\]

\noindent\underline{Note:}
This is a random function, as it depends implicitly on $X_1,\ldots,X_n$.

The \underline{log-likelihood function}
\[
\ell_n(\theta)=\sum_{i=1}^{n}\log f_\theta(X_i),
\]
we write
\[
\overline{\ell}_n(\theta)=\frac{1}{n}\ell_n(\theta)
\]
for a \underline{normalised log-likelihood}.

\paragraph{Def 4.}
A \underline{maximum likelihood estimator}
$\widehat{\theta}_{\mathrm{MLE}}$ is any estimator with
\[
L_n\bigl(\widehat{\theta}_{\mathrm{MLE}}\bigr)
=
\max_{\theta\in\Theta}L_n(\theta).
\]

\paragraph{Ex.}
\[
X_i\overset{\mathrm{iid}}{\sim}\operatorname{Poi}(\theta).
\]
\[
L_n(\theta)
=
\prod_{i=1}^{n}e^{-\theta}\frac{\theta^{X_i}}{X_i!}.
\]
\[
\ell_n(\theta)
=
-n\theta+(\log\theta)\sum X_i+\mathrm{const.}
\]
\[
\ell_n'(\theta)
=
-n+\frac{\sum X_i}{\theta}
\overset{\mathrm{FOC}}{=}0
\]
\[
\Longrightarrow\quad
\text{suggests }
\widehat{\theta}_{\mathrm{MLE}}
=
\frac{\displaystyle\sum_{i=1}^{n}X_i}{n}
:=
\overline{X}_n
\qquad
\text{``sample mean''}.
\]

\noindent\underline{SOC:}
\[
\ell_n''(\theta)
=
-\frac{1}{\theta^2}\sum X_i
<0
\qquad
\text{if }\sum X_i>0.
\]

If $\sum X_i=0$,
\[
\ell_n(\theta)=-n\theta
\quad\Longrightarrow\quad
\widehat{\theta}_{\mathrm{MLE}}=0.
\]

In every case
\[
\widehat{\theta}_{\mathrm{MLE}}=\overline{X}_n.
\]

\paragraph{Ex.}
\[
X_i\overset{\mathrm{iid}}{\sim}N(\mu,\sigma^2),
\qquad
\theta=(\mu,\sigma^2).
\]
\[
\ell_n(\theta)
=
-\frac{n}{2}\log\pi
-\frac{n}{2}\log(\sigma^2)
-\frac{1}{2\sigma^2}
\sum_{i=1}^{n}(X_i-\mu)^2.
\]

\noindent\underline{FOC:}
\[
\frac{\partial\ell_n}{\partial\mu}
=
-\frac{1}{\sigma^2}
\sum_{i=1}^{n}(X_i-\mu)
=0,
\]
\[
\frac{\partial\ell_n}{\partial\sigma^2}
=
-\frac{n}{2\sigma^2}
+
\frac{1}{2(\sigma^2)^2}
\sum_{i=1}^{n}(X_i-\mu)^2
=0.
\]

We can check that the Hessian of $\ell_n$ is negative definite
$\forall\,\theta\in\mathbb{R}\times(0,\infty)$.

Hence,
\[
\widehat{\theta}_{\mathrm{MLE}}
=
\begin{pmatrix}
\overline{X}_n \\[6pt]
\underbrace{
\displaystyle\frac{1}{n}
\sum_{i=1}^{n}(X_i-\overline{X}_n)^2
}_{\text{``Sample variance''}}
\end{pmatrix}.
\]

\subsection*{Non-examinable}

The likelihood principle poses that if two observations
\[
X=(X_1,\ldots,X_n)
\qquad\text{and}\qquad
X^*=(X_1^*,\ldots,X_n^*)
\]
give rise to likelihoods $L_n,L_n^*$ respectively, such that
\[
L_n(\theta)=L_n^*(\theta)\cdot M
\qquad \forall\,\theta\in\Theta,
\]
and some constant $M$, then they should produce the same
``inference'' on the true parameter.

\begin{itemize}
    \item The MLE, Bayes estimators, \underline{credible sets}
    satisfy the likelihood principle.
    \item \underline{Confidence sets}, tests generally do not.
\end{itemize}

\noindent\underline{Birnbaum (1962)}
\[
\text{Conditionality}+\text{Sufficiency}
\quad\Longrightarrow\quad
\text{Likelihood principle}.
\]

\begin{center}
end of lecture 1
\end{center}
\end{document}
